\section{A single error budget for continuous-state approximation}\label{sec:transfer}
A finite-state theorem alone does not justify replacing a continuous model by a grid. The following sufficient conditions permit such a transfer without restricting the original policy to be piecewise constant.

\begin{assumption}[Certified reset approximation]\label{ass:reset}
For a finite Borel partition $(E_i)_{i=1}^{n_h}$, let $\mu_i$ be probability measures supported on $E_i$ and write $\nu=\sum_i\nu_i\mu_i$. There are cellwise rewards $r_{h,t}$, costs $k_{h,t}\in[0,K]$, and terminal values $g_h$, and a kernel
\[
 Q_{h,t}^a(x,dy)=\sum_j P_{h,t}^a(i,j)\mu_j(dy),\qquad x\in E_i,
\]
with certified uniform errors $d_r,d_k,d_g$ and $\alpha_h$ satisfying
\[
 |r-r_h|\leq d_r,\quad |k-k_h|\leq d_k,\quad |g-g_h|\leq d_g,
 \quad |(P_t^a-Q_{h,t}^a)f(x)|\leq\alpha_h\norm{f}_\infty
\]
for every bounded Borel $f$. The approximation rewards obey the same bounds $R,G$ as the original rewards, or $R,G$ are enlarged to cover both models. All finite transition entries and error bounds are computable with verified directions.
\end{assumption}

The last inequality uses the operator convention for total variation; it is not the half-$\ell_1$ convention for action probabilities above. Define $M_{J,t}=R\sum_{s=0}^{T-t-1}\beta^s+G\beta^{T-t}$ and $M_{C,t}=K\sum_{s=0}^{T-t-1}\beta^s$. Set
\begin{align}
 a_T&=d_g,& a_t&=d_r+\beta a_{t+1}+\beta\alpha_hM_{J,t+1},\label{eq:primitiveJ}\\
 b_T&=0,& b_t&=d_k+\beta b_{t+1}+\beta\alpha_hM_{C,t+1},\label{eq:primitiveC}\\
 a_h&=\max_t a_t,& b_h&=b_0.\nonumber
\end{align}
Let $v_h(s)$ be the optimum of the finite common-Markov model at operating allowance $s$ and initial masses $(\nu_i)$.

\begin{theorem}[End-to-end transfer]\label{thm:transfer}
Under Assumption~\ref{ass:reset}, if $\varepsilon>2a_h$, then
\begin{equation}\label{eq:transfer}
 v_h(\varepsilon+2a_h)-b_h
 \leq\KR(\nu,\varepsilon)
 \leq v_h(\varepsilon-2a_h)+b_h.
\end{equation}
Suppose finite computation returns intervals $[L_+,U_+]$ and $[L_-,U_-]$ for these relaxed and tightened optima, of widths at most $\eta_+$ and $\eta_-$. A valid global interval is
\begin{equation}\label{eq:total}
 [\max\{0,L_+-b_h\},U_-+b_h],
\end{equation}
and its width is at most
\begin{equation}\label{eq:totalwidth}
 \eta_++\eta_-+2b_h+
 D_C\frac{4a_h}{(1-\beta)(\varepsilon-2a_h)+4a_h}.
\end{equation}
If all certified primitive errors tend to zero and both finite-solver widths tend to zero, this interval converges to the original measurable-policy optimum. Its measured width supplies a termination test for the actual economic objective.
\end{theorem}
\begin{proof}
Fix any measurable Markov policy $p$, not necessarily cellwise constant. Backward subtraction of its original and reset-model Bellman equations gives \eqref{eq:primitiveJ}--\eqref{eq:primitiveC}. Bellman comparison gives the same operating error for optimal values. In the reset model, averaging the policy within $E_i$, $\bar p_t(a\mid i)=\int p_t(a\mid x)\mu_i(dx)$, averages its value exactly: after every transition the within-cell law is again $\mu_j$, and stage rewards and costs are cellwise constant. This identity follows backward, including the common continuation. An original feasible policy has reset operating value at least the reset optimum minus $\varepsilon+2a_h$ pointwise, hence also after cell averaging; its finite-model cost differs from the original initial cost by at most $b_h$. Taking an infimum gives the lower inequality. Conversely, lift a finite policy feasible at $\varepsilon-2a_h$ by using its action probabilities throughout each cell. The reset value is constant in each cell; the two operating errors make the lifted original policy feasible, and its cost differs by at most $b_h$. This proves the upper inequality. Finally apply Theorem~\ref{thm:repair} in the finite model with target allowance $\varepsilon-2a_h$ and violation $4a_h$, then add the two finite-computation widths and the cost errors.
\end{proof}

\subsection{Where the transfer theorem applies}
On a compact domain, a transition density with computable uniform continuity in the source variable and an $L_1$-approximable density in the destination variable can furnish the reset approximation. The cell reference measures must also represent the initial law as in Assumption~\ref{ass:reset}. For example, with Lebesgue-uniform cell measures and a uniform initial law, a uniformly Lipschitz transition density gives a computable $\alpha_h=O(h)$ on a shape-regular mesh. Cellwise operating and cost approximations give $d_r,d_k,d_g=O(h)$ when their moduli are Lipschitz. Aligned discontinuities in an installed rule can be retained as exact partition boundaries. In this case \eqref{eq:totalwidth} is $O(h+\eta_++\eta_-)$ at fixed positive $\varepsilon$ and $\beta<1$, with explicitly displayed constants. The number of cells still scales as $h^{-\dim X}$; the theorem does not remove the curse of dimensionality.

The primary affine-shock model and the nonlinear two-state model below have atomic transition laws. Smearing their successor atoms into cells generally leaves the total-variation error bounded away from zero. We therefore do \emph{not} apply Theorem~\ref{thm:transfer} to those models by asserting a small mesh. Their reported intervals are direct whole-domain bounds for the unchanged atomic models. They remain valid without a convergence claim for that particular spatial representation. This distinction prevents the approximation theorem from being used as an unsupported explanation for closing an atomic-model gap.

\section{Certificates that accelerate the global problem}\label{sec:certificates}
The complete search and the transfer theorem provide the convergence argument. The inherited one-sided certificates are useful accelerators and diagnostic tools, but are not substituted for that argument.

\subsection{Deterministic outer restrictions and measurability}
Let $L_t\leq V_t\leq U_t$ be bounded Borel operating witnesses, including the terminal date. Every all-restart feasible deterministic policy must select an action in
\begin{equation}\label{eq:necessary}
 \cN_t^{L,U}(x)=\{a:L_t(x)-T_t^aU_{t+1}(x)\leq\varepsilon\}.
\end{equation}
Indeed, its operating continuation cannot exceed $V_{t+1}\leq U_{t+1}$. The set is nonempty because an operating-optimal action satisfies the test. Its graph is Borel, since it is a finite union of Borel sublevel sets. Backward minimization of $k_t+\beta P_t^aB_{t+1}$ over this set yields a Borel lower function; the lowest-index minimizer is a Borel selector. Thus the general-state recursion does not rely on an implicit measurable-selection assertion.

This is a necessary-action outer class, not the original constrained problem. A selected continuation can accumulate losses that violate later operating requirements. Equality between its integrated cost lower bound and a separately feasible policy's cost certifies $\KD(\nu,\varepsilon)$. Equality of the full functions certifies more. Tightening $L,U$ alone does not force equality; complete search is what controls the remaining structural gap.

\subsection{Randomized support floors and a connected local program}\label{sec:local}
For $\lambda\geq0$, a verified supersolution $W_t^\lambda$ of
\[
 W_T^\lambda\geq\lambda g,\qquad
 W_t^\lambda\geq\max_a\{\lambda r_t-k_t+\beta P_t^aW_{t+1}^\lambda\}
\]
implies $C_t^p\geq\lambda(V_t-\varepsilon)-W_t^\lambda$ for every feasible policy, including randomized Markov policies. Define a floor $F_t$ by the maximum of zero and finitely many such terms, using $L_t$ in place of $V_t$ when necessary. Given a valid next-date lower function $B_{t+1}$, define
\[
 \mathcal D_t(x)=\left\{p\in\Delta(A):\sum_ap_aT_t^aU_{t+1}(x)\geq L_t(x)-\varepsilon\right\}.
\]
\begin{equation}\label{eq:localLP}
 B_t(x)=\max\left\{F_t(x),\min_{p\in\mathcal D_t(x)}
 \sum_ap_a\{k_t(x,a)+\beta P_t^aB_{t+1}(x)\}\right\},\quad B_T=0.
\end{equation}
The feasible action distribution of any admissible policy belongs to the local constraint, and its continuation cost dominates $B_{t+1}$. Induction proves that \eqref{eq:localLP} is a global randomized cost lower bound. It does not prove that the local optimizer can be combined into a feasible policy or that finite support prices close a global duality gap.

For three actions and one local operating inequality, the local polytope has only pure feasible vertices and two-action boundary mixtures. Our new primary calculation evaluates these vertices with exact rational arithmetic on directed whole-cell bounds. Unlike the earlier diagnostic with a zero global floor, it uses the nonzero $F_t$ at every date and is paired with a feasible policy for the same uniform-initial randomized objective. It improves that lower endpoint in \LocalImproveCount{} of the 42 configurations. The prior adaptive local-LP theorem and all of its experimental records remain in the retained supplement; their target is still the local relaxation, not \eqref{eq:target}.

\subsection{Whole-domain randomized deployments}\label{sec:lottery}
A simple backward construction supplies feasible upper policies without changing the atomic model. Given a feasible next-date value $J_{t+1}$, write $q_0=T_t^{a_t^0}J_{t+1}$, $q_*=\max_aT_t^aJ_{t+1}$, and $b=V_t-\varepsilon$. If $q_0\geq b$, retain the installed action. Otherwise mix it with a lowest-index maximizing action with probability
\begin{equation}\label{eq:lottery}
 p_t^{\rm change}(x)=\frac{b(x)-q_0(x)}{q_*(x)-q_0(x)}.
\end{equation}
The denominator is positive wherever this expression is used. The resulting operating value is $J_t=\max\{q_0,b\}$ and is feasible at every state. In the piecewise-affine model, $J_t$ remains piecewise affine while the mixing probability can be a ratio of affine functions. Cost is therefore enclosed rather than falsely represented as an exact piecewise-affine function.

The verifier includes isolated switching points in every operating identity. Directed cost boxes integrate against the three absolutely continuous initial laws. All affine transition slopes are strictly positive, so a finite set of switching points and its finite-horizon preimages have zero probability under these laws. That fact justifies almost-everywhere cost enclosures; it does not justify ignoring knots in a restart operating constraint or using the same cost enclosure for a Dirac initial law.

\subsection{A verified direct mathematical-programming comparator}\label{sec:comparator}
The direct finite-state formulation minimizes $\sum_i\nu_iC_{0i}$ subject to \eqref{eq:operating}--\eqref{eq:cost}, $J_{ti}\geq V_{ti}-\varepsilon$, and the common policy simplexes. The constraints are bilinear in policy probabilities and continuation values. A McCormick relaxation on each probability box provides a classical global lower program. The baseline uses the same incumbent repair, contraction rule, box selection, node allowance, initial law, and all-restart constraints as interval Bellman search. The additional work is the stronger LP bound and an LP-proposed feasible candidate after repair.

Floating-point LP status is not a certificate. For a minimization relaxation $c'z$ subject to $Az\leq b$ and $\ell\leq z\leq u$, every rational vector $\lambda\geq0$ gives
\begin{equation}\label{eq:residual}
 \mathcal L(\lambda)=-\lambda'b+
 \sum_j\min\{(c+A'\lambda)_j\ell_j,(c+A'\lambda)_ju_j\}
 \leq\min c'z.
\end{equation}
This follows by adding $\lambda'(Az-b)\leq0$ and minimizing the resulting affine function over the box. Rounded nonnegative solver marginals are evaluated in \eqref{eq:residual} using exact rational arithmetic. A failed LP falls back to the interval Bellman lower bound; a floating infeasibility declaration never removes the policy box. A separate verifier rebuilds every relaxation and every search-tree partition from the primitive model.

